\chapter{多个多元均值向量的比较}
\phantomsection\label{chap:6}

\section{引言}
\phantomsection\label{sec:6.1}

第 5 章中所建立的统计思想可推广到处理多个均值向量的比较问题.在多元正态或大样本的假设下,其理论与方法比第 5 章略为复杂些.在讨论过程中,我们将经常回顾一元情形的处理方法,然后用类比法将其推广到相应的多元情形.我们所提供数学实例将有助于加深对概念的理解.

由于均值的比较往往(且应该)来源于实验设计,我们就有机会来讨论优秀的实验设计应遵循的一些基本原则.在行为研究中很有用的\mvbold{重复测量设计},以及这种设计在分析生长曲线时所需的一些修改,我们都作了明确考虑.

6.2 节先考虑成对均值向量的比较.在随后几节中,将讨论按处理水平排列的均值向量的若干种比较方法.相应的检验统计量依赖于总变差按各种来源的变差及误差的分解,这种分解就是所谓的\mvbold{多元方差分析}(简记为 MANOVA).

\section{成对比较与重复测量设计}
\phantomsection\label{sec:6.2}

\subsection{成对比较}

为了考察不同实验条件下的响应是否明显不同,人们常常将这些条件下的实验测量值记录下来.例如,对一种新药的效果或广告宣传的饱和度,可通过比较``处理''(新药与广告)前与处理后的测量值来确定.在另一些情形下,对相同或相似的实验单元可使用两种或更多种处理,比较其响应,便可评价各种处理引起的效应.

用来对两种处理进行比较或用来判断一种处理是否存在的一个合理方法,就是将两种处理指派给相同的单元(个人、商店、小块土地等),然后通过计算成对响应之差,消除单元到单元的外来变差所引起的大部分影响,从而完成对这些响应的分析.

在单个响应(即一元)的情形下,用 $X_{j1}$ 表示在第 $j$ 个试验下对处理 1 的响应(或处理前的响应),用 $X_{j2}$ 表示在第 $j$ 个试验下对处理 2 的响应(或处理后的响应).也就是说 $(X_{j1},X_{j2})$ 为对第 $j$ 个单元或对第 $j$ 对相似单元记录的测量值.通过设计,$n$ 个差值
\phantomsection\label{eq:6-1}
\begin{equation}
D_j = X_{j1} - X_{j2},\quad j = 1,2,\cdots,n \tag{6-1}
\end{equation}
应只反映各种不同处理所引起的效应.

若已知\hyperref[eq:6-1]{式(6-1)} 中的差值 $D_j$ 为来自分布 $N(\delta,\sigma_d^2)$ 的独立观测值,则变量
\phantomsection\label{eq:6-2}
\begin{equation}
t = \frac{\bar{D} - \delta}{s_d/\sqrt{n}} \tag{6-2}
\end{equation}
